La organización \textbf{Voluntarios Unidos} coordina actividades de apoyo comunitario en tres áreas: \textbf{educación (E)}, \textbf{salud (S)} y \textbf{medio ambiente (M)}. Los voluntarios tienen un tiempo limitado semanal y la organización desea maximizar el impacto de sus actividades, medido en términos de puntos de beneficio para la comunidad.

Por cada hora dedicada, cada área genera los siguientes beneficios. Educación: \textbf{5 puntos} de beneficio, salud: \textbf{8 puntos} de beneficio y medio ambiente: \textbf{6 puntos} de beneficio.

La organización cuenta con un máximo de \textbf{40 horas de trabajo voluntario semanalmente}. 

Además, por política interna, al menos el \textbf{30\% del tiempo total} debe destinarse a actividades educativas.

Por último, la organización desea un beneficio de al menos 200 puntos entre las actividades de educación y salud.

Si $x_1$, $x_2$, $x_3$ representan el tiempo semanal dedicado (en horas) a educación, salud y medio ambiente, respectivamente, el modelo de programación lineal que permite maximizar el impacto comunitario es:


\begin{equation}
\begin{split}
\text{max. } z = 5x_1 + 8x_2 + 6x_3 \\
\text{s.a.:} \\
x_1 + x_2 + x_3 \leq 40 \quad \text{(Tiempo total disponible)} \\
0.7x_1 - 0.3x_2 - 0.3x_3 \geq 0 \quad \text{(Política educativa)}\\
5x_1 + 8x_2 \geq 200 \quad \text{(Límite en salud y medio ambiente)} \\
\end{split}
\end{equation}

Que se puede formular de manera equivalente de la siguiente manera:


\begin{equation}
\begin{split}
\text{max. } z = 5x_1 + 8x_2 + 6x_3 \\
\text{s.a.:} \\
x_1 + x_2 + x_3 \leq 40 \quad \text{(Tiempo total disponible)} \\
7x_1 - 3x_2 - 3x_3 \geq 0 \quad \text{(Política educativa)} \\
5x_1 + 8x_2 \geq 200 \quad \text{(Límite en salud y medio ambiente)} \\
x_1, x_2, x_3 \geq 0 \quad \text{(No negatividad)}
\end{split}
\end{equation}

Se conoce la tabla correspondiente a la solución óptima del problema:
\begin{center}
\begin{tabular}{c|c|cccccc|}
 & $z$ & $x_1$ & $x_2$ & $x_3$ & $h_1$ & $h_2$ & $h_3$\\ 
Phase 2 & -284 & 0 & 0 & -2 & $-71/10$ & $-3/10$ & 0 \\ 
\hline
$h_3$ & 84  & 0 & 0 & 8 & $71/10$ & $3/10$ & 1\\ 
$x_1$ & 12  & 1 & 0 & 0 & $3/10$ & $-1/10$ & 0\\ 
$x_2$ & 28  & 0 & 1 & 1 & $7/10$ & $1/10$ & 0\\ 
\hline
\end{tabular}
\end{center}



Se pide:

\begin{enumerate}
    \item \textbf{Interpretar la información que ofrece la tabla de la solución óptima en términos de las actividades que se cumplen, los beneficios generados y el cumplimiento de requisitos.}

    La programación óptima de actividades consiste en:
    \begin{itemize}
        \item Dedicar 12 horas a educación, 28 horas a salud y ninguna a medio ambiente.
        \item El beneficio total obtenido sería de 284 puntos.
        \item Se haría uso del tiempo total disponible ($h_1=0$).
        \item La proporción de tiempo dedicada a actividades educativas se cumpliría por el mínimo, 30\% ($h_2=0$).
        \item Se cumple con el mínimo de beneficio de las actividades de educación y salud superando el mínimo en 84 puntos ($h_3=84$).
    \end{itemize}

    \item \textbf{Obtener la inversa de la base correspondiente a la solución óptima.}

    Las columna de la inversa de la base de obtienen de la siguiente manera.
    \begin{itemize}
        \item Primera columna: columna de las tasas de sustitución correspondiente a $h_1$ (restricción de tipo menor o igual).
        \item Segunda columna: opuesta de la columna de las tasas de sustitución correspondiente a $h_2$ (restricción de tipo mayor o igual).
        \item Tercera columna: opuesta de la columna de las tasas de sustitución correspondiente a $h_3$ (restricción de tipo mayor o igual).
    \end{itemize}

\begin{equation}
\begin{split}
B^{-1} =\begin{pmatrix}
71/10 & -3/10 & -1\\
3/10 & 1/10 & 0\\
7/10 & -1/10 & 0\\
\end{pmatrix}
\end{split}
\end{equation}
    
    \item \textbf{Evaluar el impacto sobre las horas dedicadas a las diferentes actividades y sobre el beneficio de las actividades realizadas si, en una semana en particular, fuera necesario dedicar una hora a actividades relacionadas con medio ambiente.}

    Este cambio quedaría reflejado en el modelo de programación lineal por el hecho de que $x_3$ tomaría valor 1.  

    \begin{itemize}
        \item El impacto sobre el beneficio se puede calcular a partir del criterio del simplex de $x_3$. Como $V^B_{x_3} = -2$, realizar una hora de actividades de medio ambiente daría lugar a una reducción del beneficio de 2 unidades (sería igual a 282).
        \item El impacto sobre las horas a las actividades se podría evaluar estudiando los cambios de los valores de $x_1$ y $x_2$ y estos cambios vienen dados por las tasas de sustitución correspondientes.
        \begin{itemize}
            \item $p^B_{23} = 0$ representa la disminución del valor de la segunda variable básica ($x_1$) cuando la tercera variable aumenta su valor una unidad ($x_3$=1). Se dedicaría el mismo tiempo a educación, se decir, 12 horas.
            \item $p^B_{33} = 1$ representa la disminución del valor de la tercera variable básica ($x_3$) cuando la tercera variable aumenta su valor una unidad ($x_3$=1). Se dedicaría una hora menos a salud, es decir, 27 horas.
        \end{itemize}
    \end{itemize}

    Por un lado, al dedicar una hora más a medio ambiente, tendríamos 6 puntos de beneficio adicionales y, por otro, al dedicar una hora menos en salud, tendríamos 8 puntos de beneficio menos: el resultado neto es una reducción de 2 unidades del beneficio, como indica $V^B_{x_3}$.


    \item \textbf{Cuantificar el impacto sobre las horas dedicadas a cada actividad y sobre los puntos de beneficio si, como resultado de la celebración de la semana de voluntariado se dispusiera de un número genérico de horas adicionales a las 40 semanales igual a $H$.}

    Si dispusiéramos de $H$ horas adicionales a las 40, cambiaría el término $b$ del problema, que sería:

    \begin{equation}
    b'=
    \begin{pmatrix}
    40 + H \\
    0\\
    200
    \end{pmatrix}    
    \end{equation}

    Al cambiar b, cambiarían los valores de $u^B$ y $z^B$.

    \begin{equation}
    \begin{split}
    u^B=B^{-1}b = \begin{pmatrix}
    71/10 & -3/10 & -1\\
    3/10 & 1/10 & 0\\
    7/10 & -1/10 & 0\\
    \end{pmatrix}
    \begin{pmatrix}
    40 + H\\
    0\\
    200\\
    \end{pmatrix}
     = \begin{pmatrix}
    84+H/10\\
    12+3H/10\\
    28+7H/10
    \end{pmatrix}
    \end{split}
    \end{equation}

    La solución básica correspondiente nunca dejaría de ser factible y se realizarían $3H/10$ horas más de educación y $7H/10$ de salud.

    El impacto sobre los puntos sería el siguiente:

    \begin{equation}
    \begin{split}
    z^B=c^Bu^B = \begin{pmatrix}
    0 & 5 & 8\\
    \end{pmatrix}
    \begin{pmatrix}
    84\\
    12+3H/10\\
    28+7H/10\\
    \end{pmatrix}
     = 284 + 71H/10\end{split}
    \end{equation} 

    Es decir, el beneficio se incrementaria en 71H/10.

    Todo lo anterior se puede observar igualmente, analizando qué ocurriría si hacemos que $h_1 = -H$ que sería análogo a disponer de $H$ horas adicionales de tiempo total disponible.

    Siguiendo un razonamiento análogo al de la pregunta anterior y dado el carácter lineal del problema, se podría realizar el análisis evaluando el impacto de que $\Delta H = -H$, con el mismo resultado.
    

    

    
\end{enumerate}


